\section{Trajectory Sampling and Coverage Optimization}
\label{sec:algo}

The coverage structure of $\sigma(S)$ suggests a sampling-based approach
to seed selection. Instead of explicitly computing the consumption
probabilities $M_{u,s}$, we simulate independent coupon trajectories and
record their terminal consumers. Grouping these outcomes by consumer
produces an inverted coverage index, on which greedy selection can be
performed. We refer to this procedure as \CIMRIS{} to indicate its
connection to reverse coverage representations in influence
maximization~\cite{borgs2014maximizing,tang2015influence}.
Here, reverse coverage refers to indexing candidate seeds by their
sampled consumers; the samples themselves are generated by forward
simulation, rather than by traversing incoming edges from a root.

This section establishes the sampling distribution, the coverage
formulation, and a $(1-1/e-\varepsilon)$ approximation guarantee.
Throughout, $1\le k\le n$, and we use the condition
$p_u^a+p_u^d>0$ for every $u\in V$ stated in
Section~\ref{sec:model}. Write
$\beta=\min_{u\in V}(p_u^a+p_u^d)>0$.
The parameter $\beta$ depends on the input instance.

\subsection{Independent Trajectory Samples}
\label{sec:source-worlds}

For each sample index $t$ and candidate seed $s\in V$, let
\[
  \widetilde W_s^{(t)}
  =\bigl(r_{u,h}^{(s,t)}\bigr)_{u\in V,\ h\ge1},
  \qquad r_{u,h}^{(s,t)}\sim\mathrm{Uniform}[0,1),
\]
where all variables are mutually independent across all indices
$(s,t,u,h)$. The index $h$ counts visits to node $u$ by the coupon
originating at $s$. At each visit, the coupon is consumed, dropped,
or transferred along $(u,v)$ with probabilities $p_u^a$, $p_u^d$,
and $p(u,v)$, respectively. In particular, revisiting a node uses
a fresh variable, not the action sampled at an earlier visit.
Assigning an independent world to every candidate seed is valid
because each selected seed receives one coupon and distinct coupons
propagate independently. Unselected candidates are used only to
define the common empirical objective for seed selection.

Let $Z_s^{(t)}\in V\cup\{\bot\}$ be the terminal outcome, with
$Z_s^{(t)}=u$ indicating consumption by $u$ and $Z_s^{(t)}=\bot$
indicating dropping. By the definition of $M_{u,s}$,
\begin{equation}
\label{eq:endpoint-law}
  \Pr[Z_s^{(t)}=u]=M_{u,s},\qquad
  \Pr[Z_s^{(t)}=\bot]=1-\sum_{u\in V}M_{u,s}.
\end{equation}
If $L_s^{(t)}$ counts node visits, including the terminal visit,
then $\Pr[L_s^{(t)}>h]\le(1-\beta)^h$ for every integer $h\ge0$.
Thus trajectories terminate almost surely and
$\mathbb E[L_s^{(t)}]\le1/\beta$.

Algorithm~\ref{alg:kernel} samples one terminal outcome. For an ordering
$v_1,\ldots,v_{d_u}$ of the out-neighbors of $u$, precompute
\[
 b_{u,0}=p_u^a+p_u^d,\qquad
 b_{u,i}=b_{u,0}+\sum_{\ell=1}^{i}p(u,v_\ell).
\]
Equation~\eqref{eq:normalize} gives $b_{u,d_u}=1$.
The forwarding branch uses the unconditional edge probabilities
$p(u,v)$ directly. No additional factor $p_u^t$ is applied.

\begin{algorithm}[t]
\caption{\textsc{SampleTerminalConsumer}$(s)$}
\label{alg:kernel}
\begin{algorithmic}[1]
\Require Seed $s$; node probabilities; precomputed boundaries $b_{u,i}$
\Ensure Terminal outcome in $V\cup\{\bot\}$
\State $u\gets s$
\While{true}
  \State Draw a fresh $r\sim\mathrm{Uniform}[0,1)$
  \If{$r<p_u^a$}
    \State \Return $u$
  \ElsIf{$r<p_u^a+p_u^d$}
    \State \Return $\bot$
  \Else
    \State Find $i$ such that $b_{u,i-1}\le r<b_{u,i}$
    \State $u\gets v_i$
  \EndIf
\EndWhile
\end{algorithmic}
\end{algorithm}

\subsection{Inverted Coverage Sets}
\label{sec:reverse-coverage}

\begin{definition}[Inverted coverage set]
For a sample index $t$ and consumer $u\in V$, define
\begin{equation}
\label{eq:reverse-set-def}
  R_t(u)=\{s\in V:Z_s^{(t)}=u\}.
\end{equation}
The indexed family $\mathcal R_t=(R_t(u))_{u\in V}$ contains one
set for each potential consumer.
\end{definition}

To generate $\mathcal R_t$, simulate one independent trajectory
from every candidate $s\in V$ and insert $s$ into the bucket
indexed by its terminal consumer, unless its coupon is dropped.
Each candidate belongs to at most one bucket in a family, so the
sets $R_t(u)$ are pairwise disjoint for fixed $t$.

\begin{lemma}[Coverage identity]
\label{lem:coverage}
For every fixed seed set $S\subseteq V$ and consumer $u$,
\begin{equation}
\label{eq:reverse-coverage}
 \Pr[S\cap R_t(u)\ne\emptyset]
 =1-\prod_{s\in S}(1-M_{u,s}).
\end{equation}
\end{lemma}
\begin{proof}
The intersection is empty exactly when $Z_s^{(t)}\ne u$ for
every $s\in S$. These events depend on mutually independent
candidate worlds. Equation~\eqref{eq:endpoint-law} therefore gives
$\Pr[S\cap R_t(u)=\emptyset]=\prod_{s\in S}(1-M_{u,s})$.
Taking the complement proves the identity.
\end{proof}

Generate $\theta$ independent families, retaining each indexed
sample even when its outcomes coincide with those of another sample.
Define
\begin{equation}
\label{eq:empirical-spread}
 Y_t(S)=\sum_{u\in V}\mathbf1[S\cap R_t(u)\ne\emptyset],
 \qquad
 \hat\sigma_\theta(S)=\frac1\theta\sum_{t=1}^\theta Y_t(S).
\end{equation}
Here $Y_t(S)$ counts distinct consumers, rather than the number
of consumed coupons. In particular, $0\le Y_t(S)\le |S|$.

\begin{lemma}[Unbiasedness and empirical submodularity]
\label{lem:unbiased}
For every fixed $S\subseteq V$,
$\mathbb E[\hat\sigma_\theta(S)]=\sigma(S)$.
For every fixed collection of samples, $\hat\sigma_\theta$ is a
normalized, monotone, submodular function on $V$.
\end{lemma}
\begin{proof}
Linearity of expectation and Lemma~\ref{lem:coverage} yield
\[
 \mathbb E[\hat\sigma_\theta(S)]
 =\sum_{u\in V}\left[1-\prod_{s\in S}(1-M_{u,s})\right]
 =\sigma(S).
\]
For fixed samples, consider the coverage universe
$\mathcal U=[\theta]\times V$. Candidate $s$ covers
$(t,Z_s^{(t)})$ whenever $Z_s^{(t)}\ne\bot$.
Consequently, $\theta\hat\sigma_\theta(S)$ is the size of the
union of the items covered by candidates in $S$.
It is a standard coverage function, which has all three stated
properties.
\end{proof}

Independence is required across sample families and across
candidate trajectories. The coverage indicators for different
consumers within a family need not be independent. Accordingly,
the concentration analysis below treats $Y_t(S)$ as one random
variable for each family.

\subsection{Greedy Selection on the Sampled Objective}
\label{sec:greedy-selection}

We apply greedy selection to the fixed empirical objective
$\hat\sigma_\theta$. Maintain an indicator $C_{t,u}$ for each
covered item and an integer counter $g_s$ for each candidate.
Initially, $g_s$ is the number of non-dropped outcomes for $s$.
When an item $(t,u)$ is first covered, decrement $g_s$ for every
$s\in R_t(u)$. For every unselected candidate, this maintains
the invariant
\[
 g_s=\theta\bigl(\hat\sigma_\theta(S\cup\{s\})
                       -\hat\sigma_\theta(S)\bigr).
\]
Thus the updates implement exact greedy selection on the sampled
objective; they do not approximate its marginal gains.
All greedy iterations reuse the same samples.

The sample budget uses a deterministic lower bound on
$\mathrm{OPT}=\max_{|S|=k}\sigma(S)$. Let $S_a$ be any set of
the $k$ nodes with largest $p_u^a$, and set
\begin{equation}
\label{eq:opt-lower-bound}
  \mathrm{LB}=\sum_{u\in S_a}p_u^a.
\end{equation}

\begin{lemma}[Spread lower bound]
\label{lem:lower-bound}
$\mathrm{LB}\le\sigma(S_a)\le\mathrm{OPT}$.
If $\mathrm{LB}=0$, every seed set has zero expected spread.
\end{lemma}
\begin{proof}
For $u\in S_a$, let $I_u$ indicate that the coupon seeded at $u$
is consumed on its first visit. Every such event activates a
distinct user, so the final number of activated users is at least
$\sum_{u\in S_a}I_u$ in every realization.
Taking expectations gives $\sigma(S_a)\ge\sum_{u\in S_a}p_u^a$.
The upper bound follows from optimality. Since $k\ge1$ and $S_a$
contains the largest consumption probabilities, $\mathrm{LB}=0$
implies $p_u^a=0$ for every node, making consumption impossible.
\end{proof}

For $\mathrm{LB}>0$, choose $\varepsilon\in(0,1-1/e)$ and
$\delta\in(0,1)$, and set
\begin{equation}
\label{eq:sample-budget}
 \theta=\left\lceil
 \frac{10k}{\varepsilon^2\mathrm{LB}}
 \ln\left(\frac{2\binom nk}{\delta}\right)
 \right\rceil.
\end{equation}
One may replace $\ln\binom nk$ by the upper bound $k\ln(en/k)$
when computing this budget.

\begin{algorithm}[t]
\caption{\CIMRIS{}: Trajectory Sampling and Coverage Greedy}
\label{alg:cim-ris}
\begin{algorithmic}[1]
\Require $G$ and model probabilities; $1\le k\le n$; $\varepsilon,\delta$
\Ensure A seed set $S$ of size $k$
\State Compute $S_a$ and $\mathrm{LB}$ using Eq.~\eqref{eq:opt-lower-bound}
\If{$\mathrm{LB}=0$}
  \State \Return $S_a$
\EndIf
\State Set $\theta$ using Eq.~\eqref{eq:sample-budget}
\State Precompute the transition boundaries $b_{u,i}$
\For{$t=1,\ldots,\theta$}
  \State $R_t(u)\gets\emptyset$ for all $u\in V$
  \For{each $s\in V$}
    \State $Z_s^{(t)}\gets\Call{SampleTerminalConsumer}{s}$
    \If{$Z_s^{(t)}\ne\bot$}
      \State Insert $s$ into $R_t(Z_s^{(t)})$
    \EndIf
  \EndFor
\EndFor
\State $S\gets\emptyset$; $C_{t,u}\gets\mathrm{false}$ for all $(t,u)$
\State $g_s\gets\sum_{t=1}^\theta\mathbf1[Z_s^{(t)}\ne\bot]$ for all $s$
\For{$i=1,\ldots,k$}
  \State Choose $v\in\arg\max_{s\in V\setminus S}g_s$
  \State $S\gets S\cup\{v\}$
  \For{$t=1,\ldots,\theta$}
    \State $u\gets Z_v^{(t)}$
    \If{$u\ne\bot$}
      \If{$C_{t,u}=\mathrm{false}$}
        \State $C_{t,u}\gets\mathrm{true}$
        \For{each $s\in R_t(u)$}
          \State $g_s\gets g_s-1$
        \EndFor
      \EndIf
    \EndIf
  \EndFor
\EndFor
\State \Return $S$
\end{algorithmic}
\end{algorithm}

\subsection{Approximation Guarantee}
\label{sec:theoretical-guarantees}

\begin{theorem}
\label{thm:approx}
Under the stated model assumptions, Algorithm~\ref{alg:cim-ris}
returns a set $\hat S$ of size $k$ such that, with probability
at least $1-\delta$,
\begin{equation}
\label{eq:approx-bound}
 \sigma(\hat S)\ge(1-1/e-\varepsilon)\mathrm{OPT}.
\end{equation}
\end{theorem}
\begin{proof}
The claim is immediate if $\mathrm{LB}=0$. Otherwise, fix a set
$S$ of size $k$. The variables $Y_t(S)$ are independent across
$t$, lie in $[0,k]$, and have mean $\sigma(S)$. Moreover,
\[
 \operatorname{Var}(Y_t(S))\le\mathbb E[Y_t(S)^2]
 \le k\sigma(S).
\]
Set $a=\varepsilon\mathrm{OPT}/2$. Bernstein's inequality gives
\begin{align}
\label{eq:bernstein}
 \Pr\bigl[|\hat\sigma_\theta(S)-\sigma(S)|\ge a\bigr]
 &\le 2\exp\left(-\frac{\theta a^2}
 {2k\sigma(S)+(2/3)ka}\right)\nonumber\\
 &\le 2\exp\left(-\frac{\theta\varepsilon^2\mathrm{OPT}}
 {10k}\right).
\end{align}
The last inequality uses $\sigma(S)\le\mathrm{OPT}$ and
$4(2+\varepsilon/3)<10$. Taking a union bound over all
$\binom nk$ sets and using $\mathrm{LB}\le\mathrm{OPT}$ and
Eq.~\eqref{eq:sample-budget}, the probability of any such
deviation is at most
\[
 2\binom nk\exp\left(-\frac{\theta\varepsilon^2\mathrm{LB}}
 {10k}\right)\le\delta.
\]
On the complementary event, every size-$k$ set has estimation
error at most $a$, including the data-dependent output $\hat S$.
Let $S^*$ maximize the true objective and write $\alpha=1-1/e$.
Greedy maximization of the fixed coverage objective gives
$\hat\sigma_\theta(\hat S)\ge\alpha\hat\sigma_\theta(S^*)$
by the standard cardinality-constrained submodular guarantee~\cite{nemhauser1978analysis}.
Consequently,
\begin{align*}
 \sigma(\hat S)
 &\ge\hat\sigma_\theta(\hat S)-a
 \ge\alpha\hat\sigma_\theta(S^*)-a\\
 &\ge\alpha\mathrm{OPT}-(1+\alpha)a
 \ge(\alpha-\varepsilon)\mathrm{OPT}.
\end{align*}
\end{proof}

\subsection{Computational Cost and Scope}
\label{sec:discussion-sampling}

Computing $S_a$ by sorting takes $O(n\log(n+1))$ time, and
precomputing the transition boundaries takes $O(n+m)$ time.
With binary search for forwarding decisions, the expected cost
per trajectory is $O(\log(n+1)/\beta)$, including the constant
cost of a terminal action. Generating all families therefore
takes $O(n\theta\log(n+1)/\beta)$ expected time.

There are at most $n\theta$ memberships in the inverted index.
Each bucket is processed at most once, when its item is first
covered, so all counter decrements together cost $O(n\theta)$.
Scanning candidate gains costs $O(nk)$, while examining the
selected endpoints costs $O(k\theta)\le O(n\theta)$.
Thus, when $\mathrm{LB}>0$, the total expected time is
\begin{equation}
\label{eq:total-complexity}
 O\left(m+n\log(n+1)+\frac{n\theta\log(n+1)}{\beta}+nk\right),
\end{equation}
with $O(n+m+n\theta)$ memory. When $\mathrm{LB}=0$, the algorithm
returns after computing the lower bound.

These bounds expose the cost of the construction: each family
requires $n$ forward trajectories, and the sufficient sample
budget can be large when $\mathrm{LB}$ is small or the requested
accuracy is high. The procedure avoids explicitly storing the
matrix $M$ and reuses sampled outcomes throughout greedy selection,
but these properties alone do not establish a runtime or memory
advantage over other evaluation methods.

The construction also differs from assigning one shared world
to all hypothetical starting nodes. Here, candidate-specific
independence is what makes the union-coverage identity match
independent coupons. A directly generated reverse set cannot
replace $R_t(u)$ solely because it satisfies
$\Pr[s\in R_t(u)]=M_{u,s}$ for each $s$: the joint membership
distribution must also support Eq.~\eqref{eq:reverse-coverage}.
Developing alternative backward constructions is a separate
question and is not required for the guarantee above.
